\begin{solapartat}
La superfície de una sola espira és:
\[ s = 0,025^2 = 6,25\cdot 10^{-4}\,\mathrm{m^2}\ \text{\punts{0,1}} \]
El flux del camp magnètic que travessa la bobina, l'escriurem com:
% DUBTE: a l'original «cos 0·» (probablement 0°); es transcriu tal com és.
\[ \Phi = \vec{B}\cdot\vec{S} = B\cdot S\cdot\cos 0\cdot = B\cdot S\ \text{\punts{0,2}} \]
on $S$ serà:
\[ S = 2000\cdot s = 1,25\ \mathrm{m^2}\ \text{\punts{0,1}} \]
A partir de la lectura de la gràfica podem escriure:
\[ B(t)_{t\in[0, 5]} = \frac{25 - 0}{5 - 0}\cdot 10^{-3}\cdot t\ \mathrm{T}\ \text{\punts{0,1}} \Rightarrow \Phi(t)_{t\in[0, 5]} = 6,25\cdot 10^{-3}\cdot t\ \mathrm{Wb}\ \text{\punts{0,2}} \]
\[ B(t)_{t\in[5, 8]} = 25\cdot 10^{-3}\ \mathrm{T}\ \text{\punts{0,1}} \Rightarrow \Phi(t)_{t\in[5, 8]} = 31,3\cdot 10^{-3}\ \mathrm{Wb}\ \text{\punts{0,2}} \]
\end{solapartat}

\begin{solapartat}
\[ \varepsilon(t) = -\frac{\mathrm{d}\Phi}{\mathrm{d}t} = -\frac{\mathrm{d}B}{\mathrm{d}t}\cdot S\ \text{\punts{0,4}} \]
\[ \varepsilon(t)_{t\in[0, 5]} = -6,3\cdot 10^{-3}\ \mathrm{V}\ \text{\punts{0,2}} \]
\[ \varepsilon(t)_{t\in[5, 8]} = 0\ \mathrm{V}\ \text{\punts{0,2}} \]
\[ \varepsilon(t)_{t\in[8, 10]} = 16\cdot 10^{-3}\ \mathrm{V}\ \text{\punts{0,2}} \]
Es considerarà igualment correcte si la $\varepsilon((t)$ del primer tram es positiva i la del últim es negativa. També es considerarà igualment correcte si es realitza el càlcul sense la utilització del càlcul diferencial i es duu a terme considerant al quocient de les variacions:
\[ \varepsilon = -\frac{\Delta\Phi}{\Delta\mathrm{t}} \]
\end{solapartat}
